\documentclass[10pt,a4paper]{amsart}
\usepackage[margin=19mm]{geometry}
\usepackage{amsmath,amssymb,mathtools}
\usepackage[scr=rsfs,cal=euler]{mathalfa}
\usepackage{xfrac}
\usepackage[unicode,hidelinks,bookmarks=false]{hyperref}
\newcommand{\ie}{i.e.\ }
\newcommand{\cf}{cf.\ }
\newcommand{\resp}{respectively\ }
\newcommand{\Subsection}{\subsection}
\providecommand{\nobreakdash}{\nobreak\hbox{-}}
\setlength{\emergencystretch}{2em}
\multlinegap0pt
\usepackage{algorithm}\usepackage{algpseudocode}
\usepackage{float}
\usepackage{amssymb}
\newtheorem{theorem}{Theorem}
\title{On the Condition Number Upper Bound of the L-BFGS Inverse Hessian Approximation Matrix with a Two-Sided Geometric Envelope Safeguarding Mechanism}
\author{Don Li}
\date{Source: arXiv:2607.05836v1 (CC BY 4.0)}
\begin{document}
\maketitle

\noindent\emph{Source and licence.} On the Condition Number Upper Bound of the L-BFGS Inverse Hessian Approximation Matrix with a Two-Sided Geometric Envelope Safeguarding Mechanism. Version arXiv:2607.05836v1 (CC BY 4.0), distributed by arXiv under the Creative Commons Attribution 4.0 International licence, http://creativecommons.org/licenses/by/4.0/, as recorded in the arXiv licence metadata for this exact version and shown on the abstract page https://arxiv.org/abs/2607.05836v1. Full licence notice: \texttt{licenses/arxiv-2607.05836-CC-BY-4.0.txt}.\par\smallskip

\noindent\textit{Editorial scope.} Counted unit: the article's theorem theorem (source0.tex lines 429--437). The article's complete original proof of it is delivered with it (source0.tex lines 439--611, one \texttt{proof} environment of 173 lines). No mathematics is written, repaired or re-derived: the statement and the proof are the article's own lines, in the article's own notation, and no retained line is substituted or re-flowed.  No further source-local context is needed: every cross-reference of every retained line resolves inside the two delivered spans.  Nothing was omitted from any retained span, so the package carries no omission record.  No retained line contains a citation, so the wrapper carries no bibliography and no bibliography entry.  No retained line contains a figure environment, an image-inclusion command or drawing code, so no image is delivered and the package declares no asset dependency.  Editorial formatting only, in addition: an AMS \texttt{amsart} preamble that restates the article's own declarations and macros used by the retained lines, namely the environments \texttt{theorem} on separate counters, together with the standard packages those lines need.  The retained environments print in this document as Theorem 1 in order of appearance, whereas the article prints the same items with the numbering of its own full document.  The complete native source of the article as distributed in the arXiv e-print for this exact version is bundled unchanged as \texttt{sources/204576/source0.tex}.
\begin{theorem}[Uniform Upper Bound on Two-Sided L-BFGS $\kappa(H_{k+1})$]
Let \(\{H_{k+1}\}_{k=0}^{\infty}\) be the sequence of inverse Hessian update approximation matrices generated by the Two-Sided L-BFGS algorithm (cf. Algorithm 3) with memory depth $m$, Li-Fukushima lower bound $\epsilon > 0$, and upper bound $M \gg 0$. Then, for any objective function \(f: \mathbb{R}^n \to \mathbb{R}\) and any iteration $k$, the condition number \(\kappa(H_{k+1})\) is uniformly bounded above by a finite constant \(\kappa _{\max }\) depending solely on the dimension \(n\), the memory depth \(m\), and the envelope parameters \(\epsilon \) and \(M\), i.e., 

\begin{equation} \kappa(H_{k+1}) \leq \kappa_{\max} < \infty, \quad \forall k \geq 0, 
\end{equation}

\noindent
where \(\kappa_{\max} = \frac{M_1}{m_1}\), and \(M_1, m_1\) are the positive, finite roots of the transcendental equation \(\lambda - \ln \lambda = C^*-n+1\), for any eigenvalue $\lambda$ of $H_{k+1}$, derived from the uniform upper bound \(C^*(n, m, \epsilon, M)\) of the Trace-Determinant lemma.
\end{theorem}

\begin{proof}
As in L-BFGS with standard OS scaling, at any iteration $k$, Two-Sided L-BFGS constructs $H_{k+1}$ via at most $m$ successive BFGS updates to the scaled identity matrix, i.e., $H_{k+1}^{(0)} = \gamma_{k+1}I$. By construction, Two-Sided L-BFGS only saves vector pairs $(s_k, y_k)$ at iteration $k$ to the memory buffer iff $\frac{y_k^T s_k}{\|s_k\|^2} \geq \epsilon$ for $\epsilon > 0$ and $\frac{\|y_k\|^2}{y_k^T s_k}$ for $M \gg 0$. Let the memory buffer of length $m$ be indexed by $j = 1,2, ..., \hat{m}$, where $\hat{m} \leq m$. Then the sequence of matrix updates in Two-Sided L-BFGS is

\begin{equation}
    \begin{cases}
        H_{k+1}^{(0)} = \gamma_{k+1}I \quad \text{for }j=0, \\
        H_{k+1}^{j} = (I - \rho_k s_jy_j^T)H_{k+1}^{(j-1)}(I - \rho_j y_j s_j^T) + \rho_j s_j s_j^T, \quad \text{for } j \in [1, \hat{m}].
    \end{cases}
\end{equation}

\noindent
The final matrix is $H_{k+1} = H_{k+1}^{(\hat{m})}$. By the Trace-Determinant lemma, we have

\begin{equation}
\psi(H_{k+1}^{(j)}) \leq \psi(H_{k+1}^{(j-1)}) + \frac{\|s_j\|^2}{y_j^T s_j} - 1 - \text{ln}(\frac{y_j^T s_j}{s_j^T B_{k+1}^{(j-1)}s_j}),
\end{equation}

\noindent
where $B_{k+1}^{(j-1)} = [H_{k+1}^{(j-1)}]^{-1}$. Summing inequality (20) telescopically for $j$ up to $\hat{m}$ yields

\begin{equation}
\psi(H_{k+1}) \leq \psi(\gamma_{k+1}I) + \sum_{j=1}^{\hat{m}}\frac{\|s_j\|^2}{y_j^T s_j} - \hat{m} - \sum_{j=1}^{\hat{m}}\text{ln}(\frac{y_j^T s_j}{s_j^T B_{k+1}^{(j-1)}s_j}).
\end{equation}

\noindent
Thus, 

\begin{equation}
\psi(H_{k+1}) \leq \psi(\gamma_{k+1}I) + \sum_{j=1}^{\hat{m}}\frac{\|s_j\|^2}{y_j^T s_j} - \sum_{j=1}^{\hat{m}}\text{ln}(\frac{y_j^T s_j}{s_j^T B_{k+1}^{(j-1)}s_j}),
\end{equation}

\noindent
or equivalently,

\begin{equation}
\psi(H_{k+1}) \leq \psi(\gamma_{k+1}I) + \sum_{j=1}^{\hat{m}}\frac{\|s_j\|^2}{y_j^T s_j} + \sum_{j=1}^{\hat{m}}\text{ln}(\frac{s_j^T B_{k+1}^{(j-1)}s_j}{y_j^T s_j}).
\end{equation}

\noindent
To upper bound $\psi(H_{k+1})$, it suffices to upper bound each of the three terms on the right-hand side of (23), respectively. \\
\\
We first find an upper bound on $\psi(\gamma_{k+1}I)$. Recall $\gamma_{k+1} = \frac{s_k^Ty_k}{\|y_k\|^2}$. By the ``upper side" of the envelope condition, $\frac{||y_k||^2}{y_k^Ts_k} \leq M$. Then $\gamma_{k+1} \geq \frac{1}{M}$. Likewise, the Li-Fukushima cautious update rule acts as the ``lower side" of the envelope condition, where $\frac{y_k^T s_k}{\|s_k\|^2} \geq \epsilon$. By the Cauchy-Schwarz inequality, $y_k^T s_k \leq \|y_k\| \|s_k\|$, which in turn implies $\|s_k\| \geq \frac{y_k^T s_k}{\|y_k\|}$. Substituting this inequality into the Li-Fukushima bound yields $\epsilon \leq \frac{y_k^T s_k}{\|s_k\|^2} \leq \frac{(y_k^T s_k)}{(\frac{y_k^T s_k}{\|y_k\|})^2} = \frac{\|y_k\|^2}{y_k^T s_k}$. Note that the Li-Fukushima bound also implies $\|s_k\| \leq \frac{\|y_k\|}{\epsilon}$. Then we have $\gamma_{k+1} = \frac{s_k y_k^T}{\|y_k\|^2} \leq \frac{\|s_k\| \|y_k\|}{\|y_k\|^2} = \frac{\|s_k\|}{\|y_k\|} \leq \frac{(\frac{\|y_k\|}{\epsilon})}{\|y_k\|} = \frac{1}{\epsilon}$, so $\gamma_{k+1} \leq \frac{1}{\epsilon}$. Therefore, the OS factor itself in Two-Sided L-BFGS is constrained such that $\gamma_{k+1} \in [\frac{1}{M}, \frac{1}{\epsilon}]$, so $\frac{1}{M} \leq \gamma_{k+1} \leq \frac{1}{\epsilon}$. By the Trace-Determinant lemma, $\psi(\gamma_{k+1}I) = \text{tr}(\gamma_{k+1}I) - \text{ln}[\text{det}(\gamma_{k+1})I] = n\gamma_{k+1} - n\text{ln}(\gamma_{k+1})$. Since the function $h(\gamma) = n\gamma - n\text{ln}(\gamma)$ is convex and continuous over $[\frac{1}{M}, \frac{1}{\epsilon}]$, $h(\gamma)$ is bounded above at the endpoints of $[\frac{1}{M}, \frac{1}{\epsilon}]$. Therefore,

\begin{equation}
\psi(\gamma_{k+1}I) \leq \text{max}\{\frac{n}{M}-n\text{ln}(\frac{1}{M}), \frac{n}{\epsilon}-n\text{ln}(\frac{1}{\epsilon})\} < \infty.
\end{equation}

\noindent
We next upper bound $\sum_{j=1}^{\hat{m}} \frac{\|s_j\|^2}{y_j^T s_j}$. By the Li-Fukushima bound, $\frac{\|s_j\|^2}{y_j^T s_j} \leq \frac{1}{\epsilon}$, so we have

\begin{equation}
\sum_{j=1}^{\hat{m}} \frac{\|s_j\|^2}{y_j^T s_j} \leq \frac{\hat{m}}{\epsilon} \leq \frac{m}{\epsilon} < \infty,
\end{equation}

\noindent
since $\hat{m} \leq m$. Lastly, we upper bound $\sum_{j=1}^{\hat{m}}\text{ln}(\frac{s_j^T B_{k+1}^{(j-1)}s_j}{y_j^T s_j})$. Consider $\frac{s_j^T B_{k+1}^{(j-1)}s_j}{y_j^T s_j}$. Note $\frac{s_j^T B_{k+1}^{(j-1)}s_j}{y_j^T s_j} = (\frac{s_j^T B_{k+1}^{(j-1)}s_j}{\|s_j\|^2})(\frac{\|s_j\|^2}{y_j^T s_j})$, so we can upper bound each of $\frac{s_j^T B_{k+1}^{(j-1)}s_j}{\|s_j\|^2}$ and $\frac{\|s_j\|^2}{y_j^T s_j}$, respectively. We already showed that $\frac{\|s_j\|^2}{y_j^T s_j} \leq \frac{1}{\epsilon}$ immediately follows from the Li-Fukushima bound. Since $B_{k+1}^{(j-1)}$ is symmetric, it follows from its Rayleigh quotient bound that $\frac{s_j^T B_{k+1}^{(j-1)}s_j}{\|s_j\|^2} \leq \lambda_{\text{max}}(B_{k+1}^{(j-1)})$. Now we want to upper bound $\lambda_{\text{max}}(B_{k+1}^{(j-1)})$. The update rule for $B_{k+1}^{(j-1)}$ is

\[
B^{(j)} = B^{(j-1)} - \frac{B^{(j-1)}s_j s_j^T B^{(j-1)}}{s_j^T B^{(j-1)}s_j} + \frac{y_j y_j^T}{y_j^T s_j}.
\]

\noindent
Taking the trace (Tr) of both sides yields

\[
\text{Tr}(B^{(j)}) = \text{Tr}(B^{(j-1)}) - \frac{\|B^{(j-1)}s_j\|^2}{s_j^T B^{(j-1)}s_j} + \frac{\|y_j\|^2}{y_j^T s_j}.
\]

\noindent
Since $- \frac{\|B^{(j-1)}s_j\|^2}{s_j^T B^{(j-1)}s_j} < 0$, this also gives

\[
\text{Tr}(B^{(j)}) \leq \text{Tr}(B^{(j-1)}) + \frac{\|y_j\|^2}{y_j^T s_j}.
\]

\noindent
But by the ``upper side" of the geometric envelope of Two-Sided L-BFGS, $\frac{\|y_j\|^2}{y_j^T s_j} \leq M$, so we then have

\[
\text{Tr}(B^{(j)}) \leq \text{Tr}(B^{(j-1)}) + M.
\]

\noindent
This inequality represents a telescoping sum from $j = 1$ to $\hat{m}$. For initialization $B^{(0)} = \frac{1}{\gamma_{k+1}}I$, this telescoping sum evaluates to

\[
\text{Tr}(B^{(\hat{m})}) \leq \text{Tr}(\frac{1}{\gamma_{k+1}}I) + \hat{m}M.
\]

\noindent
But recall we showed that $\gamma_{k+1} \geq \frac{1}{M}$, which implies $\frac{1}{\gamma_{k+1}} \leq M$, so $\text{Tr}(\frac{1}{\gamma_{k+1}}I) \leq nM$. Moreover, since $\hat{m} \leq m$, this gives us

\[
\text{Tr}(B^{(\hat{m})}) \leq nM + mM = (n+m)M.
\]

\noindent
Since $B^{(\hat{m})}$ is symmetric positive definite with $j \leq \hat{m}$, $\lambda_{\text{max}}(B_{k+1}^{(j-1)}) \leq \text{Tr}(B^{(\hat{m})})$. Thus,

\[
\lambda_{\text{max}}(B_{k+1}^{(j-1)}) \leq (n+m)M,
\]

\noindent
so $\frac{s_j^T B_{k+1}^{(j-1)}s_j}{\|s_j\|^2} \leq (n+m)M$. Then we yield

\[
\sum_{j=1}^{\hat{m}} \text{ln}(\frac{s_j^T B_{k+1}^{(j-1)}s_j}{y_j^T s_j}) \leq \sum_{j=1}^{\hat{m}} \text{ln}((n+m)M)(\frac{1}{\epsilon}) = \sum_{j=1}^{\hat{m}}\text{ln}(\frac{(n+m)M}{\epsilon}) = \hat{m}\text{ln}(\frac{(n+m)M}{\epsilon}) \leq m \text{ln}(\frac{(n+m)M}{\epsilon}).
\]

\noindent
Thus, we yield the upper bound

\begin{equation}
\sum_{j=1}^{\hat{m}}\text{ln}(\frac{s_j^T B_{k+1}^{(j-1)}s_j}{y_j^T s_j}) \leq m\text{ln}(\frac{(n+m)M}{\epsilon}) < \infty.
\end{equation}

\noindent
Then summing inequalities (24), (25), and (26) gives

\begin{equation}
\psi(H_{k+1}) \leq \text{max}\{\frac{n}{M}-n\text{ln}(\frac{1}{M}), \frac{n}{\epsilon}-n\text{ln}(\frac{1}{\epsilon})\} + \frac{m}{\epsilon} + m\text{ln}(\frac{(n+m)M}{\epsilon}) < \infty.
\end{equation}

\noindent
Let $C^{\star}(n,m, \epsilon, M) := \text{max}\{\frac{n}{M}-n\text{ln}(\frac{1}{M}), \frac{n}{\epsilon}-n\text{ln}(\frac{1}{\epsilon})\} + \frac{m}{\epsilon} + m\text{ln}(\frac{(n+m)M}{\epsilon})$, so we likewise have

\begin{equation}
\psi(H_{k+1}) \leq C^{\star}(n,m, \epsilon, M) < \infty.
\end{equation}

\noindent
Byrd and Nocedal's potential function for $H_{k+1}$ can likewise be written as

\begin{equation}
\psi(H_{k+1}) = \sum_{i=1}^{n}(\lambda_i - \text{ln}\lambda_i)
\end{equation}

\noindent
for eigenvalues $\lambda_i$ of $H_{k+1}$. For any arbitrary eigenvalue $\lambda_k$ of $H_{k+1}$, we have

\begin{equation}
\psi(H_{k+1}) = (\lambda_k - \text{ln}\lambda_k) + \sum_{i \neq k}^{n}(\lambda_i - \text{ln}\lambda_i).
\end{equation}

\noindent
$h(\lambda) = \lambda - \text{ln}\lambda$ has a global minimum at $\lambda = 1$, and $h(1) = 1$. Thus, $\sum_{i \neq k}^{n}(\lambda_i - \text{ln}\lambda_i) \geq (n-1)(1) = n-1$. This inequality and Eqn. (28) together imply

\begin{equation}
\lambda_k - \text{ln}\lambda_k \leq \psi(H_{k+1}) - n + 1.
\end{equation}

\noindent
Since $\psi(H_{k+1}) \leq C^\star$, we equivalently have

\begin{equation}
\lambda_k - \text{ln}\lambda_k \leq C^\star - n + 1.
\end{equation}

\noindent
Because the function $h(\lambda) = \lambda - \ln\lambda$ satisfies $\lim_{\lambda \to 0^+} h(\lambda) = \infty$ and $\lim_{\lambda \to \infty} h(\lambda) = \infty$, the transcendental equation $\lambda - \ln\lambda = C^* - n + 1$ possesses exactly two positive roots, $m_1$ and $M_1$ where $m_1 \leq M_1$. Therefore, for all eigenvalues $\lambda_i$ of $H_{k+1}$, $\lambda_i(H_{k+1}) \in [m_1, M_1]$. Thus,

\begin{equation}
0 < m_1 \leq \lambda_i(H_{k+1}) \leq M_1 < \infty, \quad \forall i \in [1,n].
\end{equation}

\noindent
Hence, $\kappa(H_{k+1})$ is uniformly bounded by

\begin{equation}
\kappa(H_{k+1}) = \frac{\lambda_{\text{max}}(H_{k+1})}{\lambda_{\text{min}}(H_{k+1})} \leq \frac{M_1}{m_1} = \kappa_{\text{max}}(H_{k+1}) < \infty.
\end{equation}
\end{proof}

\end{document}
